\section{Input Structure}
\label{sec:input-format}

FEMaster reads a unified keyword-based input language. A command
begins with an asterisk, optional parameters follow after commas,
and keyword-specific data are provided on subsequent lines.
Supported Abaqus and CalculiX input forms coexist with FEMaster
extensions in the same deck; selecting a particular file format
or parser is not required.

Familiar Parts, Assemblies, Steps, loads, and boundary conditions
can therefore be combined with FEMaster features such as named
load and support collectors, explicit LOADCASE definitions,
generic Fields, and solver controls. Each keyword page presents
a preferred example and, where meaningful, additional supported
alternatives below it rather than in parallel columns.

\subsection{Command lines}

A command line begins with \texttt{*}. The command name follows immediately and may be followed by comma-separated parameters. Parameters can be key/value pairs or flags. Data lines continue until the next command line begins.

\begin{syntax}
*COMMAND, KEY=value, FLAG
value, value, value
value, value, value
\end{syntax}

Command and parameter names are case insensitive. Whitespace inside command names is ignored, so compact names such as \cmd{SOLIDSECTION} and spaced names such as \cmd{SOLID SECTION} refer to the same command where that keyword is supported. The manual normally uses the spaced form when it improves readability and the compact form where it is customary in native FEMaster decks.

Parameter values retain their meaning. Numerical values are read as integer or floating-point quantities according to the keyword definition. Boolean-style options may appear as flags or as explicit values where documented. Names are used to reference Parts, Instances, sets, surfaces, materials, sections, coordinate systems, Fields, amplitudes, collectors, and other model entities.

\subsection{Data lines and multiline records}

The number and meaning of data-line entries depend on the keyword. Some commands accept one record per line, such as \cmd{NODE}. Others read a fixed-size group of values, a matrix, a list of identifiers, or multiple records. The keyword pages describe the required ordering explicitly.

Comma-separated values preserve the visibility of optional empty entries. Whitespace around values is ignored. Several native FEMaster commands use the literal \val{NAN} to mean that a generalized component is not prescribed. This is different from zero: zero is a physical numerical value, whereas an empty or NaN component may mean ``leave this degree of freedom unspecified'' for the corresponding keyword.

Long connectivity or data records may continue over more than one physical line where that command permits multiline input. This representation accommodates higher-order elements and large fixed-size material or Section definitions.

\subsection{Scopes}

Several keywords open scopes. A Part contains reusable mesh topology and part-level regions. An Assembly contains Instance placement and assembly-level region definitions. Material subcommands apply to the active material. Native FEMaster \cmd{LOADCASE} opens an analysis definition. Abaqus \cmd{STEP} opens an Abaqus analysis step.

Scopes are defined by keywords, not by indentation. Indentation may be used for readability but has no syntactic meaning. Closing commands such as \cmd{END PART}, \cmd{END ASSEMBLY}, \cmd{END INSTANCE}, \cmd{END STEP}, and native \cmd{END} make the end of the corresponding scope explicit.

Parts, Assemblies, and Instances are optional for simple native FEMaster models. Nodes and elements may be defined directly without constructing an explicit Assembly. The structured form represents reusable Parts, repeated Instances, independent local ID spaces, and assembly-level selections.

\paragraph{Optional native \cmd{MODEL} marker.}
Native FEMaster may optionally place \cmd{MODEL} at root level as a descriptive marker. It does not open a scope, create a second model object, or select among several models; omitting it has no effect on the topology or analysis. \key{NAME} is optional and may be used to attach a descriptive model name.

\begin{syntax}
*MODEL, NAME=BRACKET_MODEL
\end{syntax}

This marker is part of the native input-file framing rather than a separate model-definition hierarchy level. Abaqus input does not require or use a corresponding wrapper.

\subsection{Names, identifiers, and references}

Node and element IDs are user-defined non-negative integer labels. They may start at \textbf{0}; they do not need to start at 1 and do not need to be contiguous. IDs only need to be unique within the topology scope in which they are defined. For example, \val{0}, \val{10}, and \val{1000} are valid node IDs.

Inside separate Parts, the same numerical ID may be reused because each Part has its own local topology. Instances preserve these local labels from the user's point of view. When an assembly-level command needs to identify an entity from a particular Instance, the documented Instance-qualified reference syntax is used.

Named resources are used throughout the rest of the deck. Node sets, element sets, surfaces, materials, profiles, orientations, Fields, amplitudes, collectors, and other definitions are referenced by name. Stable names expose the model structure without repeatedly addressing individual numerical IDs.

\emph{User IDs} and \emph{local positional indices} represent different concepts. User node and element IDs may be sparse and may include zero. By contrast, indices such as a Field's \val{local_node}, \val{local_ip}, or \val{local_mp} identify a position inside an element and are explicitly zero based where documented.

\subsection{FEMaster and Abaqus examples in this manual}

Each keyword page shows a primary definition with useful alternatives
stacked directly below it. A given keyword may have Abaqus-compatible
or FEMaster-specific data forms; their placement and limitations are
defined by that keyword rather than by separate input dialects.
For example, \cmd{SOLVER} belongs inside a \cmd{LOADCASE} or inside
a \cmd{STEP} after its procedure definition.

The comparison shows how the same model operation is expressed in either accepted input style without exposing parser implementation details.

\keywordpage{HEADING}

\cmd{HEADING} adds descriptive text to an input deck. It does not change the finite-element equations or the numerical result. The keyword keeps exported, archived, and regression decks self-describing and is accepted in both native FEMaster and Abaqus input.

A heading is metadata conventionally placed near the beginning of a file. Later model commands do not reference it by name, and it does not substitute for a Part name, material name, loadcase name, or other actual model identifier.

Native FEMaster also accepts the optional root-level \cmd{MODEL, NAME=...} marker described above. \cmd{HEADING} is free descriptive text, whereas \cmd{MODEL} is a compact native file marker.

\syntaxheading

\begin{syntax}
*HEADING
Descriptive model title
\end{syntax}

\exampleheading

\begin{inputdeck}[Abaqus-compatible]
*HEADING
Cantilever verification model
\end{inputdeck}
